\section{Conclusions}

The sensitivity analyses conducted in this document confirm the robustness of
the main report's principal findings across four dimensions:

\begin{itemize}
  \item \textbf{Variable selection and model diagnostics.} Among the four
    candidate specifications tested, the RatesVolEquity model (\(\Delta VIX\),
    \(\Delta S\&P\), \(\Delta UST10\), \(\Delta USYC\), \(\Delta SOFR\))
    achieves the highest out-of-sample \(R^2\) and the lowest AIC/BIC,
    with all coefficients statistically significant at the 5\% level. The
    remaining models either lack explanatory power out of sample
    (InflationMacro, FXBasisPolicy) or rely on variables that are partially
    tautological with the target (CrossMarket). These results reinforce the
    selection of the five-factor rates--volatility--equity specification in the
    main report.

  \item \textbf{Lag structure and cross-correlations.} The ACF and PACF of
    CDX IG daily changes reveal only a mild mean-reversion pattern at lag~1,
    with no significant autocorrelation beyond the first few lags.
    Cross-correlation functions against all candidate drivers are
    overwhelmingly concentrated at lag~0, confirming that a contemporaneous
    specification is appropriate. The only notable exception is a modest
    lag-1 signal from European credit indices (ITRX IG, ITRX XO), attributable
    to the time-zone difference between European and US market closes rather
    than a genuine predictive lead.

  \item \textbf{Regime-switching fair value (RV vs VIX).} Conditioning static
    OLS on a simple two-state rule based on S\&P realised volatility (from
    daily index returns) versus VIX
    does not improve out-of-sample fit for the best candidate model on this
    sample (see generated metrics in the regime section), though regime timing is
    illustrated in Figure~\ref{fig:regime_rv_vix}.
    This remains a sensitivity check; the main report continues to use the
    standard pooled static specification.

  \item \textbf{Rolling-window fair value model.} The static
    global-coefficient model achieves higher explanatory power
    (\(R^2 \approx 0.70\)) than the 180-day rolling alternative
    (\(R^2 \approx 0.50\)) over this sample, indicating that the
    spread--driver relationships are sufficiently stable for a fixed-parameter
    approach. Nevertheless, the rolling specification adapts more rapidly during
    regime transitions (March 2020, mid-2022) and provides a useful real-time
    complement for monitoring fair-value dislocations.
\end{itemize}

Taken together, these results support the analytical framework adopted in the
main report. The key modelling choices---contemporaneous regressors, a
five-factor rates--volatility--equity feature set, and a static estimation
approach---are each validated by the corresponding sensitivity test, while the
rolling-window and lag analyses provide useful extensions for practitioners
seeking adaptive or forward-looking applications.
